\documentclass[10pt,a4paper]{amsart}
\usepackage[margin=19mm]{geometry}
\usepackage{amsmath,amssymb,mathtools}
\usepackage[scr=rsfs,cal=euler]{mathalfa}
\usepackage{xfrac}
\usepackage[unicode,hidelinks,bookmarks=false]{hyperref}
\newcommand{\ie}{i.e.\ }
\newcommand{\cf}{cf.\ }
\newcommand{\resp}{respectively\ }
\newcommand{\Subsection}{\subsection}
\providecommand{\nobreakdash}{\nobreak\hbox{-}}
\setlength{\emergencystretch}{2em}
\multlinegap0pt
\usepackage{bm}
\usepackage{bbm}
\usepackage{dsfont}
\usepackage{xspace}
\usepackage{cancel}
\usepackage{mathtools}
\usepackage{amsthm}
\makeatletter
\@ifundefined{theorem}{\newtheorem{theorem}{Theorem}}{}
\@ifundefined{remark}{\newtheorem{remark}[theorem]{Remark}}{}
\makeatother
\def\R{\mathbb{R}} % Reals
\renewcommand{\epsilon}{\varepsilon} % Nice looking epsilon
\title[Power-Saving Bounds For Monic Minkowski Polynomials]{Power-Saving Bounds For Monic Minkowski Polynomials}
\author{Seamus Lavine}
\date{Source: arXiv:2606.30690v1 (CC BY 4.0)}
\begin{document}
\maketitle

\noindent\emph{Source and licence.} Power-Saving Bounds For Monic Minkowski Polynomials. Version arXiv:2606.30690v1 (CC BY 4.0), distributed under the Creative Commons Attribution 4.0 International licence, as recorded in the arXiv licence metadata for this exact version and shown on the abstract page https://arxiv.org/abs/2606.30690v1. Full rights notice: licenses/arxiv-2606.30690-CC-BY-4.0.txt.\par\smallskip

\noindent\textit{Editorial scope.} Counted unit: the Remark whose source label is \texttt{None} in the article (source0.tex lines 261--263, source label \texttt{None}), delivered together with the article's own complete original proof of it (source0.tex lines 265--335, one \texttt{proof} environment of 71 lines). No mathematics is written, repaired or re-derived: statement and proof are the article's own text line for line. Required context retained uncounted: none. Every definition, display and label of those lines is retained unchanged. Omitted from this package and not redistributed: 1--260, 264--264, 336--389. Nothing inside a retained excerpt is omitted or altered: every retained body is delivered exactly as the article's lines are written, so no omission receipt of any kind accompanies this package. Citations: the retained excerpts carry no citation command at all, so no bibliography entry of the article is transcribed and no reference is redistributed. Reference adaptation: none was needed and none was made; every reference inside the delivered unit resolves inside it, so no unresolved reference or citation is produced. Printed slips kept literal: no printed word, symbol, hypothesis or number of the counted statement or of its proof was altered. Layout and class: the article's own class is \texttt{amsart} with packages including amsmath, amssymb, amsthm, amsfonts, mathtools, bm, xcolor, graphicx, xy, tikz, tikz-cd; the wrapper is the lane's pinned \texttt{amsart} editorial wrapper at 10pt on A4, which restates the theorem-like environments the retained text needs and adds the article's own macro definitions that the retained text needs (\texttt{\textbackslash R}, \texttt{\textbackslash epsilon}), with extra wrapper packages (bm, bbm, dsfont, xspace, cancel, mathtools). The delivered numbering is the unit's own. Assets: no retained span contains a figure environment, a graphics-inclusion command, a PDF-page inclusion, a table or a TikZ picture; no image file is copied into this package, the package declares no asset dependency and no image file is redistributed with it. Licence: version arXiv:2606.30690v1 of this article is distributed under the Creative Commons Attribution 4.0 International licence, as recorded in the arXiv licence metadata for this exact version and shown on the abstract page https://arxiv.org/abs/2606.30690v1. A full rights notice is delivered at licenses/arxiv-2606.30690-CC-BY-4.0.txt. Characters: every retained excerpt is pure ASCII, so the delivered engine, XeLaTeX with its default Latin Modern text font, reports no missing character.
\begin{remark}
    In what follows, the constants $C_0,C_1,C_2$ are absolute and given by the lemmas. The constant $C_3$ depends on $f$, and subsequent constants depend on $f$ and the previous constants. 
\end{remark}

\begin{proof}[Proof of Theorem 1.1] For arbitrarily large $d$, let $K$ be a totally real number field of degree $d$ with embeddings $\sigma_1,\dots,\sigma_d:K\hookrightarrow\R$. By Theorem 2.3, we can take the discriminant $\Delta_K\leq C_2^d$ for an absolute $C_2>0$, independent of $d$. Let $Y$ be a sufficiently large constant, and we will choose a sufficiently large integer $X$ below; both $X$ and $Y$ are independent of $d$. Let 
\[
G=B^\times(Y)\quad\text{and}\quad P=X+B^+(\epsilon X),
\]
and set 
\[
\mathcal A=GP\subset \mathcal O_K.
\]
By Lemma 2.1, we have
\[
C_1X^dY^{d-1}\Delta_K^{-3/2}\leq |\mathcal A|\leq X^dY^dC_1^{-d}.
\tag{1}
\]
If $f(x)=\sum_{j=0}^k a_jx^j$ with $a_k=1$ and, without loss of generality, $a_0=0$, we can write
\[
f(\mathcal A)=\bigcup_{u_k\in G^k}\bigcup_{\substack{u_{j,i}\in G^j\\1\leq j\leq k-1,\;1\leq i\leq |a_j|}}\left(u_kP^k+\sum_{j=1}^{k-1}\operatorname{sgn}(a_j)\sum_{i=1}^{|a_j|}u_{j,i}P^j\right),
\]
using the convention $\operatorname{sgn}(0)=0$. Now fix $u_k\in G^k$ and $u_{j,i}\in G^j$, $(1\leq j\leq k-1,\;1\leq i\leq |a_j|)$, and let
\[
s=u_k q_k+\sum_{j=1}^{k-1}\operatorname{sgn}(a_j)\sum_{i=1}^{|a_j|}u_{j,i}q_{j,i},
\]
where $q_k\in P^k$ and $q_{j,i}\in P^j$. Then
\[
u_k^{-1}s=q_k+\sum_{j=1}^{k-1}\operatorname{sgn}(a_j)\sum_{i=1}^{|a_j|}u_k^{-1}u_{j,i}q_{j,i},
\]
and for every fixed embedding $\sigma_\ell$, we have
\[
|\sigma_\ell(u_k^{-1}s)|\leq|\sigma_\ell(q_k)|+\sum_{j=1}^{k-1}\sum_{i=1}^{|a_j|}|\sigma_\ell(u_k^{-1}u_{j,i})||\sigma_\ell(q_{j,i})|.
\]
Since $q_k,q_{j,i}$ lie in the corresponding $P$-boxes and
$u_k,u_{j,i}$ lie in the corresponding $G$-boxes, we have
\[
|\sigma_\ell(q_k)|\leq (1+\epsilon)^kX^k,\qquad
|\sigma_\ell(q_{j,i})|\leq (1+\epsilon)^jX^j,\qquad
|\sigma_\ell(u_k^{-1}u_{j,i})|\leq e^{2kY}.
\]
By choosing $X\geq (1+\epsilon)^{k-1}e^{2kY}$, we have
\[
|\sigma_\ell(u_k^{-1}s)|\leq(1+\epsilon)^kX^k+\sum_{j=1}^{k-1}|a_j|X^k.
\]
Assuming \(\epsilon\leq 1\), we have
\[u_kP^k+\sum_{j=1}^{k-1}\operatorname{sgn}(a_j)\sum_{i=1}^{|a_j|}u_{j,i}P^j\subset u_kB^+(C_3X^k),
\]
where $C_3=2^k+\sum_{j=1}^{k-1}|a_j|$. Thus
\[
\bigcup_{\substack{u_{j,i}\in G^j\\1\leq j\leq k-1,\;1\leq i\leq |a_j|}}\left(u_kP^k+\sum_{j=1}^{k-1}\operatorname{sgn}(a_j)\sum_{i=1}^{|a_j|}u_{j,i}P^j\right)
\subset u_kB^+(C_3X^k).
\]
Since multiplication by a unit is a bijection on $\mathcal O_K$, we have 
\[
|u_kB^+(C_3X^k)|=|B^+(C_3X^k)|\leq (2C_3X^k+1)^d\leq C_4^d X^{kd},
\]
for a constant $C_4>0$, where the last inequality follows by Lemma 2.2. Since $G^k\subset B^\times(kY)$, by Lemma 2.2 we have $|G^k|\leq(C_5Y)^{d-1}$ for a constant $C_5>0$. Thus, 
\[
|f(\mathcal A)|\leq|G^k|C_4^dX^{kd}\leq C_6^dY^{d-1}X^{kd}
\tag{2}
\]
for a constant $C_6>0$. Now combining $(1)$ and $(2)$, we have
\[
\frac{|f(\mathcal A)|}{|\mathcal A|^k}\leq \frac{Y^{k-1}}{C_1^k}\left(\frac{C_7}{Y^{k-1}}\right)^d
\]
for a constant $C_7>0$. By choosing $Y$ sufficiently large so that $C_7/Y^{k-1}<1$ and absorbing the extra $Y^{k-1}/C_1^k$ term into the exponent for sufficiently large $d$, we see that there is an $\alpha\in(C_7/Y^{k-1},1)$ with
\[
|f(\mathcal A)|\leq \alpha^d|\mathcal A|^k.
\]
To conclude, the upper bound in $(1)$ gives $|\mathcal A|\leq(XYC_1^{-1})^d$. Thus by choosing $c=\frac{-\log\alpha}{\log(XYC_1^{-1})}$, we have
\[
|f(\mathcal A)|\leq|\mathcal A|^{k-c}.
\]
Since $\sigma_1$ is an injective field homomorphism, taking $A=\sigma_1(\mathcal A)$ completes the proof.
\end{proof}

\end{document}
